\chapter{Introducción general}
\label{chap:intro}

\section{Motivación}

Las crisis financieras de las últimas dos décadas han dejado una regularidad difícil de ignorar: cuando un sistema bancario frágil coincide con un deterioro agudo de las perspectivas de crecimiento, el costo de financiamiento de los gobiernos se dispara de forma más que proporcional a lo que cualquiera de los dos factores explicaría por separado. La literatura académica, sin embargo, ha tratado estos dos canales —la fragilidad bancaria sistémica y el riesgo a la baja del crecimiento— de manera predominantemente independiente. Por un lado, la literatura del nexo banca-soberano ha documentado extensamente cómo la fragilidad del sistema bancario se transmite al riesgo soberano a través de pasivos contingentes y rescates fiscales \citep{FarhiTirole2018b,AcharyaDrechslerSchnabl2014b}. Por otro, la literatura de \textit{Growth-at-Risk} ha caracterizado con precisión creciente cómo las condiciones financieras determinan de forma asimétrica la cola izquierda de la distribución del crecimiento \citep{ABG2019b}. Ninguna de las dos tradiciones ha articulado sistemáticamente ambos canales como codeterminantes \textit{conjuntos y no lineales} del riesgo soberano.

Esta tesis nace de la conjetura de que esa articulación no es un simple ejercicio de agregar un regresor: la fragilidad bancaria y el riesgo de cola del crecimiento se \textbf{complementan}, no se suman, en su transmisión al spread soberano. La fragilidad del sistema bancario se comporta como un pasivo contingente cuyo costo esperado para el fisco depende críticamente del estado del ciclo económico; cuando el crecimiento ya está en su escenario adverso, el mismo nivel de fragilidad bancaria implica un rescate fiscal proporcionalmente más costoso, y por tanto una prima de riesgo soberano desproporcionadamente mayor.

\section{El objetivo de esta investigación}

Esta investigación retoma y pone a prueba, a escala global y con un \textit{pipeline} de datos propio, la hipótesis con la que se aprobó su tema en diciembre de 2025 (Informe de Taller de Tesis I, profesor guía Juan Francisco Martínez Sepúlveda). Ese informe fijó como objetivo general \textit{evaluar empíricamente el impacto conjunto, estructural y no lineal de la fragilidad del sistema bancario ($JLoss$) y el riesgo de cola del crecimiento económico ($GaR$) sobre el spread de crédito soberano en una muestra de economías emergentes}, bajo la hipótesis central de que \textbf{el mercado penaliza de forma no lineal la coexistencia de fragilidad bancaria y vulnerabilidad macroeconómica}. Sus cuatro objetivos específicos —construir y validar $JLoss$ y $GaR$ para una muestra amplia, asegurando su reproducibilidad; cuantificar el efecto marginal de $JLoss$ sobre el spread soberano ($\beta_1$); cuantificar el efecto marginal de $GaR$ ($\beta_2$); y estimar y testear la complementariedad de ambos riesgos mediante el coeficiente de interacción ($\beta_3$)— organizan directamente el Capítulo~\ref{chap:paper2} de esta tesis.

La validación preliminar de esa hipótesis, sobre una muestra piloto de cuatro países (Chile, Brasil, Colombia y Perú, 2008--2022) y una especificación simple, ya insinuaba el patrón que el Capítulo~\ref{chap:paper2} confirma y precisa a escala mayor: en una regresión \textit{pooled} sin efectos fijos, $JLoss$ y $GaR$ eran individualmente significativos con el signo esperado, pero su interacción no lo era ($p=0{,}641$); en un panel con efectos fijos de país y tiempo, ocurría lo contrario —los niveles perdían significancia y la interacción emergía altamente significativa y con el signo esperado ($p<0{,}0001$)—, lo que el informe original describió como \emph{la evidencia central de la tesis}. Ese mismo informe comprometía una hoja de ruta de tres pilares para consolidar el hallazgo: replicar $JLoss$ y $GaR$ a escala global con un motor propio en vez de datos provistos por terceros; ampliar el vector de controles e implementar una estrategia de identificación más robusta frente a la endogeneidad por causalidad inversa; y someter el resultado a un análisis de robustez de submuestras de crisis frente a normalidad. Los tres compromisos se cumplieron —el segundo de una forma distinta, pero no menos rigurosa, de la originalmente anticipada—, como detalla el Capítulo~\ref{chap:paper2}.

El \textbf{Capítulo~\ref{chap:paper2}} desarrolla esa evidencia central hasta sus últimas consecuencias: estima, sobre un único panel de trece economías emergentes con la fragilidad bancaria construida desde la terminal Bloomberg mediante el motor de datos propio que el informe original comprometía como trabajo futuro, el término de interacción entre la fragilidad bancaria sistémica —medida mediante la métrica \textit{Joint Loss} ($JLoss$) de \citet{Chari2024b}— y el riesgo a la baja del crecimiento —medido mediante \textit{Growth-at-Risk} ($GaR$), siguiendo el marco de \textit{vulnerable growth} de \citet{ABG2019b}—, sobre el spread soberano medido, siguiendo a \citet{Chari2024b}, por el índice EMBI Global Diversified de J.P.\ Morgan. El esqueleto de la evidencia es una batería de regresiones estructurada al modo de \citet{Chari2024b} —cuatro modelos anidados por tres estructuras de efectos fijos por dos muestras—. El capítulo establece con solidez los canales de \textbf{nivel} —la fragilidad bancaria y el riesgo de cola del crecimiento elevan cada uno el spread soberano, en las doce especificaciones— y muestra que su \textbf{complementariedad}, aunque del signo predicho por el informe original y corroborada por un modelo de umbral, tiene una significancia estadística \textbf{condicional}: no se identifica sobre el panel completo de trece países —a diferencia de lo que sugería, con solo cuatro países y una especificación simple, la muestra piloto—, pero sí (i)~en el núcleo de economías emergentes de financiamiento externo y (ii)~al excluir los trimestres de crisis financiera aguda, en los que los respaldos oficiales desacoplan el canal doméstico. Esa condicionalidad es estructural, no un accidente estadístico ni una retractación del hallazgo preliminar: es una versión más precisa, y más difícil de refutar, de la misma hipótesis —el multiplicador de \textit{doom loop} se traslada al \emph{precio} de la deuda soberana cuando el mercado espera que el pasivo contingente bancario recaiga sobre un soberano no respaldado cuya deuda se fija al margen por inversores sensibles al riesgo—.

\section{Contribución y hoja de ruta}

En conjunto, esta tesis hace cuatro aportes.

\textbf{Primero, un encuadre.} La literatura del nexo banca--soberano y la de \textit{growth-at-risk} han evolucionado en paralelo: la primera trata el crecimiento como un fundamento de fondo, la segunda detiene su análisis en el producto sin proyectarlo sobre la prima soberana. Esta tesis las articula como un solo mecanismo —el multiplicador de \textit{doom loop} visto desde dos lados— y lo formaliza como un término de interacción no lineal $JLoss\times D$ (con $D\equiv-GaR$, el riesgo de cola del crecimiento) sobre el spread soberano. No es agregar un regresor: es sostener que los dos canales que la literatura mantiene separados son la misma cadena.

\textbf{Segundo, un hallazgo empírico con contenido económico.} La complementariedad $JLoss\times D$ ($D=-GaR$) sobre el spread es \textbf{condicional} en dos dimensiones que apuntan al mismo mecanismo: es significativa en el núcleo de economías emergentes de financiamiento externo ($\hat\beta_3=+0{,}47$, $p=0{,}023$) y al excluir los trimestres de crisis financiera aguda ($\hat\beta_3\approx+1{,}2$, $p<0{,}01$); combinadas, $\hat\beta_3=+0{,}94$ ($p=0{,}004$). El coeficiente se reporta sobre $D=-GaR$, de modo que $\hat\beta_3>0$ indica amplificación (en la parametrización con $GaR$, $\theta=-\beta_3<0$; esta es exactamente la traducción de signo de la predicción original del informe de Taller de Tesis I, que formulaba la hipótesis bajo la convención $GaR$ con $\beta_3<0$: ambas son la misma afirmación económica). El multiplicador de \textit{doom loop} se traslada al \emph{precio} de la deuda cuando el tenedor marginal —un inversor de cartera extranjero— reprecia el pasivo contingente bancario y el soberano \emph{no} está respaldado: no ocurre ni donde la deuda la mantienen inversores domésticos de horizonte largo, ni en las crisis con líneas de \textit{swap} y financiamiento oficial. El resultado no es que el efecto sea frágil, sino que es una \textbf{propiedad de a quién y en qué régimen se le fija el precio a la deuda soberana}.

\textbf{Tercero, dos hallazgos sobre la medición.} (i)~En la submuestra donde el EMBI y el CDS soberano coexisten, ambas variables dependientes arrojan el \emph{mismo} coeficiente de interacción —de modo que las diferencias entre versiones del panel provienen de la composición de la muestra, no de la métrica del spread—. (ii)~La cota de la malla de pérdidas del motor de $JLoss$ satura el valor en riesgo al 99\% en el 98\% de las observaciones, de modo que la métrica, tal como está calibrada en la literatura, mide la pérdida esperada más una contribución de cola evaluada a un nivel de estrés fijo, y debe interpretarse de forma ordinal.

\textbf{Cuarto, un \textit{pipeline} reproducible.} Toda la métrica de fragilidad se construye con un solo protocolo de extracción de Bloomberg para 113 bancos, versionado y trazable a una única fuente de verdad documentada en el repositorio.

El hilo que une los cuatro aportes es una disciplina metodológica: reportar con precisión dónde el mecanismo está establecido y dónde no. En un área donde abundan los resultados marginales presentados como robustos, delimitar exactamente el alcance de la evidencia —el efecto condicional que sí se identifica, y la pregunta sobre si la concentración bancaria lo amplifica, que la agenda futura del Capítulo~\ref{chap:paper2} (Sección~7.4) deja abierta por falta de potencia estadística— es en sí mismo una contribución a la credibilidad de la literatura, no solo a su contenido.

El resto de la tesis se organiza así: el Capítulo~\ref{chap:paper2} desarrolla la evidencia empírica en detalle; el Capítulo~\ref{chap:discusion} sintetiza sus implicancias de política y traza la agenda de investigación futura. El Anexo~\ref{chap:anexoA} tabula la procedencia de todas las series de datos.
